% =============================================================================
\section{Conclusiones}
% =============================================================================

\subsection{Recomendación principal}

La simulación de 30 réplicas por escenario, emparejadas mediante números
aleatorios comunes, proporciona evidencia estadística abrumadora ($p = 1{,}12 \times 10^{-11}$,
$t = 10{,}8$, IC de la diferencia $[11{,}03;\;16{,}18]$ pp) de que incorporar
un tercer empleado durante las horas pico reduce el porcentaje de clientes
con espera superior a 5 minutos de \textbf{16{,}7\,\%} (Escenario A) a
\textbf{3{,}1\,\%} (Escenario B), una mejora del 81{,}4\,\%.

Se \textbf{recomienda adoptar el Escenario B} para la tasa de llegadas actual
($\lambda_{\text{pico}} = 0{,}30$ clientes/min).

\subsection{Hallazgos secundarios}

\begin{itemize}
  \item El tiempo de espera medio se reduce de 2{,}19 min a 0{,}66 min.
  \item La longitud media de la cola decrece de 0{,}48 a 0{,}15 clientes.
  \item La utilización del personal baja del 57{,}2\,\% al 50{,}2\,\%,
    manteniéndose en niveles saludables que permiten absorber picos
    transitorios sin colapsar el sistema.
  \item El análisis de sensibilidad revela que la ventaja del Escenario B
    se mantiene para $\lambda_{\text{pico}} \lesssim 0{,}42$; por encima de
    ese umbral sería necesario incorporar un cuarto empleado.
\end{itemize}

\subsection{Limitaciones del modelo}

\begin{enumerate}
  \item \textbf{Disciplina de cola:} se asume FIFO estricto; no se modela
    abandono (reneging) ni impaciencia de los clientes.
  \item \textbf{Homogeneidad de empleados:} todos los servidores tienen la
    misma distribución de servicio; no se considera variabilidad inter-empleado.
  \item \textbf{Jornada determinista:} el horario de apertura y cierre, y las
    fronteras de hora pico, son fijas.
    En la práctica pueden fluctuar.
  \item \textbf{Sin descansos ni ausencias:} el modelo asume disponibilidad
    continua de los empleados.
  \item \textbf{Modelo no estacionario aproximado:} la generación de llegadas
    usa tasas constantes por tramos, no un proceso de Poisson no homogéneo
    exacto.
\end{enumerate}

\subsection{Posibles extensiones}

\begin{itemize}
  \item Incorporar abandono de cola: modelar que los clientes se van si
    la espera estimada supera un umbral.
  \item Optimizar el número óptimo de empleados extra como función de
    $\lambda_{\text{pico}}$, generalizando el análisis de sensibilidad.
  \item Ajustar las distribuciones de servicio a datos reales del local
    (análisis de bondad de ajuste).
  \item Extender el análisis de parada usando métodos secuenciales
    (semi-anchura fija) para garantizar el criterio del 5\,\% de precisión
    relativa~\citep{law2015simulation}.
\end{itemize}
